\begin{titlepage}
\vspace*{9mm}
{\sffamily\fontsize{10}{13}\selectfont\color{Accent}
 QUANTUM E POST-QUANTUM COMPUTING\par}
\vspace{18mm}
{\sffamily\fontsize{32}{37}\selectfont\color{Accent}
 Geometria,\par simulazione\par e calcolo\par}
\vspace{8mm}
{\sffamily\fontsize{15}{19}\selectfont Guida al laboratorio\par}
\vspace{3mm}
{\sffamily\fontsize{10}{13}\selectfont\color{Ink!65}\GuidaSelezione\par}
\vspace{12mm}
\begin{tikzpicture}[x=1cm,y=1cm,>=Stealth]
 \draw[Ink!20] (-.25,0)--(6.8,0);
 \draw[Ink!20] (0,-.3)--(0,3.5);
 \draw[Secondary!30,dashed] (0,0) circle[radius=2.5];
 \draw[Secondary,line width=1.3pt,->] (0,0)--(2.17,1.25);
 \draw[Accent,line width=1.3pt,->] (0,0)--(1.25,2.17);
 \draw[Accent!65,->] (1.3,0) arc[start angle=0,end angle=60,radius=1.3];
 \node[anchor=west,font=\sffamily\small,text=Secondary] at (2.3,1.25) {$\lvert\psi\rangle$};
 \node[anchor=west,font=\sffamily\small,text=Accent] at (1.4,2.35) {$U\lvert\psi\rangle$};
\end{tikzpicture}
\vfill
{\fontsize{11}{14}\selectfont Marzio De Corato · Danilo Bruschi\par}
{\sffamily\small\color{Ink!65}Università degli Studi di Milano\par}
\vspace{5mm}
{\color{Accent!35}\hrule height .5pt}
\vspace{3mm}
{\sffamily\footnotesize\color{Ink!60}Edizione ArsClassica 01 · Classica\par}
\end{titlepage}
\section*{Come leggere la guida}
\phantomsection\addcontentsline{toc}{section}{Come leggere la guida}
Il percorso accompagna le dimostrazioni interattive del corso. Le figure e i
calcoli conservano le configurazioni iniziali del laboratorio. Ogni sezione
introduce il modello, descrive la costruzione geometrica e segue gli esempi
svolti fino all’interpretazione dei risultati. Quindici esercizi guidati per
modulo seguono le rispettive sezioni teoriche. Ogni traccia richiama i
prerequisiti e conduce a uno svolgimento completo in pagina successiva;
si suggerisce di tentare il percorso prima di consultarlo.

Le identità simboliche sono esatte. I valori ottenuti dai calcoli sono
arrotondati; il simbolo \(\approx\) segnala le approssimazioni. Gli indici
distinguono componenti e registri; gli esponenti indicano potenze. Le matrici
mantengono la disposizione per righe e colonne.

I riferimenti alle slide permettono di affiancare la guida al materiale del
seminario. Il percorso locale indicato all’inizio di ciascuna sezione conduce
alla corrispondente simulazione, nella quale si possono modificare i parametri
e seguire l’evoluzione temporale.

Le espansioni più estese e le tabelle complete si trovano nelle appendici del
rispettivo capitolo. I collegamenti interni consentono di passare dalla
dimostrazione all’appendice e di tornare al punto di partenza.

\vspace{7mm}
{\sffamily\color{Accent}Il percorso dei quattro moduli\par}
\begin{enumerate}[leftmargin=6mm,itemsep=4pt]
 \item Dal determinismo al calcolo quantistico: dati iniziali, probabilità,
       interferenza e ampiezze.
 \item Geometria del calcolo quantistico: spazi, prodotti interni, operatori,
       stati composti e misura.
 \item Costruire programmi quantistici: porte, circuiti, preparazioni e
       lettura degli stati intermedi.
 \item Algoritmi, Fourier e fattorizzazione: oracoli, ricerca, stima di fase,
       aritmetica reversibile, RSA e Shor.
\end{enumerate}
\clearpage
